\documentclass[aspectratio=169,10pt]{beamer}
\usepackage[T1]{fontenc}
\usepackage[scaled=0.9]{helvet}   % stand-in for Amazon Ember (not installed)
\renewcommand{\familydefault}{\sfdefault}
\usepackage{booktabs,tabularx,multirow,makecell,colortbl,xcolor}
\usepackage{graphicx}
\usepackage{pgfplots}
\pgfplotsset{compat=1.18}
\usepackage{listings}
\usepackage{fontawesome5}

% ---- AWS template palette (dark) ----------------------------------
\definecolor{bgdark}{HTML}{161D26}
\definecolor{panel}{HTML}{232F3E}
\definecolor{awsblue}{HTML}{007DFF}
\definecolor{awspurple}{HTML}{952EFF}
\definecolor{awsmagenta}{HTML}{E433FF}
\definecolor{awsteal}{HTML}{00D6C7}
\definecolor{bodytext}{HTML}{F2F4F7}
\definecolor{inkmuted}{HTML}{A3ACB9}
\definecolor{footc}{HTML}{8B95A3}
\definecolor{rulec}{HTML}{5A6573}
\definecolor{gridc}{HTML}{333E4C}
\definecolor{vendorc}{HTML}{7B8594}
\definecolor{linkc}{HTML}{5FB0FF}
\definecolor{sectag}{HTML}{E433FF}
% extra categorical colours for charts, tuned for the dark background
\definecolor{cat1}{HTML}{EB6834}   % orange
\definecolor{cat2}{HTML}{3B8BEB}   % blue
\definecolor{cat3}{HTML}{1BAF7A}   % aqua

% footer text; override with \renewcommand{\deckfooter}{...} after \documentclass[aspectratio=169,10pt]{beamer}
\usepackage[T1]{fontenc}
\usepackage[scaled=0.9]{helvet}   % stand-in for Amazon Ember (not installed)
\renewcommand{\familydefault}{\sfdefault}
\usepackage{booktabs,tabularx,multirow,makecell,colortbl,xcolor}
\usepackage{graphicx}
\usepackage{pgfplots}
\pgfplotsset{compat=1.18}
\usepackage{listings}
\usepackage{fontawesome5}

% ---- AWS template palette (dark) ----------------------------------
\definecolor{bgdark}{HTML}{161D26}
\definecolor{panel}{HTML}{232F3E}
\definecolor{awsblue}{HTML}{007DFF}
\definecolor{awspurple}{HTML}{952EFF}
\definecolor{awsmagenta}{HTML}{E433FF}
\definecolor{awsteal}{HTML}{00D6C7}
\definecolor{bodytext}{HTML}{F2F4F7}
\definecolor{inkmuted}{HTML}{A3ACB9}
\definecolor{footc}{HTML}{8B95A3}
\definecolor{rulec}{HTML}{5A6573}
\definecolor{gridc}{HTML}{333E4C}
\definecolor{vendorc}{HTML}{7B8594}
\definecolor{linkc}{HTML}{5FB0FF}
\definecolor{sectag}{HTML}{E433FF}
% extra categorical colours for charts, tuned for the dark background
\definecolor{cat1}{HTML}{EB6834}   % orange
\definecolor{cat2}{HTML}{3B8BEB}   % blue
\definecolor{cat3}{HTML}{1BAF7A}   % aqua

% footer text; override with \renewcommand{\deckfooter}{...} after \documentclass[aspectratio=169,10pt]{beamer}
\usepackage[T1]{fontenc}
\usepackage[scaled=0.9]{helvet}   % stand-in for Amazon Ember (not installed)
\renewcommand{\familydefault}{\sfdefault}
\usepackage{booktabs,tabularx,multirow,makecell,colortbl,xcolor}
\usepackage{graphicx}
\usepackage{pgfplots}
\pgfplotsset{compat=1.18}
\usepackage{listings}
\usepackage{fontawesome5}

% ---- AWS template palette (dark) ----------------------------------
\definecolor{bgdark}{HTML}{161D26}
\definecolor{panel}{HTML}{232F3E}
\definecolor{awsblue}{HTML}{007DFF}
\definecolor{awspurple}{HTML}{952EFF}
\definecolor{awsmagenta}{HTML}{E433FF}
\definecolor{awsteal}{HTML}{00D6C7}
\definecolor{bodytext}{HTML}{F2F4F7}
\definecolor{inkmuted}{HTML}{A3ACB9}
\definecolor{footc}{HTML}{8B95A3}
\definecolor{rulec}{HTML}{5A6573}
\definecolor{gridc}{HTML}{333E4C}
\definecolor{vendorc}{HTML}{7B8594}
\definecolor{linkc}{HTML}{5FB0FF}
\definecolor{sectag}{HTML}{E433FF}
% extra categorical colours for charts, tuned for the dark background
\definecolor{cat1}{HTML}{EB6834}   % orange
\definecolor{cat2}{HTML}{3B8BEB}   % blue
\definecolor{cat3}{HTML}{1BAF7A}   % aqua

% footer text; override with \renewcommand{\deckfooter}{...} after \input{preamble}
\newcommand{\deckfooter}{\textcopyright{} \the\year, Amazon Web Services, Inc. or its affiliates.}

% ---- beamer look ---------------------------------------------------
\setbeamercolor{background canvas}{bg=bgdark}
\setbeamercolor{normal text}{fg=bodytext,bg=bgdark}
\setbeamercolor{structure}{fg=awsblue}
\setbeamercolor{frametitle}{fg=white}
\setbeamercolor{itemize item}{fg=awsteal}
\setbeamercolor{itemize subitem}{fg=inkmuted}
\setbeamercolor{enumerate item}{fg=awsteal}
\setbeamercolor{alerted text}{fg=awsmagenta}
\setbeamertemplate{navigation symbols}{}
\setbeamertemplate{itemize items}[square]
\setbeamertemplate{enumerate items}[default]
\setbeamertemplate{caption}[numbered]
\setbeamerfont{itemize/enumerate subbody}{size=\scriptsize}
\setbeamersize{text margin left=8mm,text margin right=8mm}

% section tag shown above each frame title (no divider slides)
\gdef\secname{}
\newcommand{\sect}[1]{\section{#1}\gdef\secname{#1}}
\setbeamertemplate{frametitle}{%
  \nointerlineskip\vskip3.2mm
  \ifx\secname\empty\else{\fontsize{6.5}{7}\selectfont\bfseries\color{sectag}\MakeUppercase{\secname}}\par\vskip0.9mm\fi
  {\fontsize{12.5}{14.5}\selectfont\bfseries\color{white}\insertframetitle}\par\vskip0.6mm}
\setbeamertemplate{headline}{}
\setbeamertemplate{footline}{%
  \hbox to\paperwidth{\hspace*{8mm}{\fontsize{5}{6}\selectfont\color{footc}\deckfooter}\hfill{\fontsize{6}{7}\selectfont\color{footc}\insertframenumber}\hspace*{8mm}}\vskip2.5mm}

% ---- deck macros ---------------------------------------------------
\newcommand{\hl}[1]{\textcolor{awsteal}{\textbf{#1}}}
\newcommand{\src}[1]{\vfill{\tiny\textcolor{footc}{Source: #1}}}
\newcommand{\flabel}[1]{\label{#1}\hypertarget{#1}{}}
% clickable slide references
\newcommand{\sref}[1]{\hyperlink{#1}{\textcolor{linkc}{slide~\pageref*{#1}}}}
\newcommand{\srefs}[2]{\textcolor{linkc}{slides~\hyperlink{#1}{\pageref*{#1}}--\hyperlink{#2}{\pageref*{#2}}}}
\newcommand{\plainen}[1]{\colorbox{panel}{\parbox{\dimexpr\linewidth-2\fboxsep}{\raggedright\footnotesize\textbf{In plain English:} #1}}\par\vspace{1mm}}
\newcommand{\idcard}[1]{\colorbox{panel}{\parbox{\dimexpr\textwidth-2\fboxsep}{\scriptsize #1}}}
\newcommand{\grp}[2]{\multicolumn{#1}{@{}l}{\textcolor{awsteal}{\textit{#2}}}}
\newcommand{\agi}[4]{\hyperlink{#1}{{\large\bfseries\color{awsteal}#2}\quad{\large\bfseries\color{white}#3}}\\[0.5mm]{\small\color{inkmuted}#4}\par\vspace{4mm}}
\newcommand{\bgframe}[1]{\usebackgroundtemplate{\includegraphics[width=\paperwidth,height=\paperheight]{#1}}}

\arrayrulecolor{rulec}
\lstset{basicstyle=\ttfamily\tiny\color{bodytext},breaklines=true,columns=fullflexible,
  backgroundcolor=\color{panel},frame=single,rulecolor=\color{rulec},
  keywordstyle=\color{awsteal}\bfseries,commentstyle=\color{inkmuted},showstringspaces=false}
\renewcommand{\arraystretch}{1.1}
\usetikzlibrary{arrows.meta,positioning,shapes.geometric,calc,fit,backgrounds,shadows}
\tikzset{box/.style={draw=awsblue,fill=panel,rounded corners=2pt,align=center,font=\scriptsize,inner sep=3pt,text=bodytext},
         arr/.style={-{Stealth[length=2mm]},thick,awsblue}}
\pgfplotsset{every axis/.append style={axis line style={rulec},tick style={rulec},
  tick label style={text=bodytext},label style={text=bodytext}}}

% numbered table / figure labels (one index across the deck)
\newcounter{tabno}\newcounter{figno}
\newcommand{\tabcap}[1]{\refstepcounter{tabno}{\scriptsize\bfseries\color{awsteal}Table~\thetabno\hspace{0.6em}{\color{bodytext}#1}}\par\vspace{0.2mm}}
\newcommand{\figcap}[1]{\refstepcounter{figno}{\scriptsize\bfseries\color{awsteal}Figure~\thefigno\hspace{0.6em}{\color{bodytext}#1}}\par\vspace{0.2mm}}

\newcommand{\deckfooter}{\textcopyright{} \the\year, Amazon Web Services, Inc. or its affiliates.}

% ---- beamer look ---------------------------------------------------
\setbeamercolor{background canvas}{bg=bgdark}
\setbeamercolor{normal text}{fg=bodytext,bg=bgdark}
\setbeamercolor{structure}{fg=awsblue}
\setbeamercolor{frametitle}{fg=white}
\setbeamercolor{itemize item}{fg=awsteal}
\setbeamercolor{itemize subitem}{fg=inkmuted}
\setbeamercolor{enumerate item}{fg=awsteal}
\setbeamercolor{alerted text}{fg=awsmagenta}
\setbeamertemplate{navigation symbols}{}
\setbeamertemplate{itemize items}[square]
\setbeamertemplate{enumerate items}[default]
\setbeamertemplate{caption}[numbered]
\setbeamerfont{itemize/enumerate subbody}{size=\scriptsize}
\setbeamersize{text margin left=8mm,text margin right=8mm}

% section tag shown above each frame title (no divider slides)
\gdef\secname{}
\newcommand{\sect}[1]{\section{#1}\gdef\secname{#1}}
\setbeamertemplate{frametitle}{%
  \nointerlineskip\vskip3.2mm
  \ifx\secname\empty\else{\fontsize{6.5}{7}\selectfont\bfseries\color{sectag}\MakeUppercase{\secname}}\par\vskip0.9mm\fi
  {\fontsize{12.5}{14.5}\selectfont\bfseries\color{white}\insertframetitle}\par\vskip0.6mm}
\setbeamertemplate{headline}{}
\setbeamertemplate{footline}{%
  \hbox to\paperwidth{\hspace*{8mm}{\fontsize{5}{6}\selectfont\color{footc}\deckfooter}\hfill{\fontsize{6}{7}\selectfont\color{footc}\insertframenumber}\hspace*{8mm}}\vskip2.5mm}

% ---- deck macros ---------------------------------------------------
\newcommand{\hl}[1]{\textcolor{awsteal}{\textbf{#1}}}
\newcommand{\src}[1]{\vfill{\tiny\textcolor{footc}{Source: #1}}}
\newcommand{\flabel}[1]{\label{#1}\hypertarget{#1}{}}
% clickable slide references
\newcommand{\sref}[1]{\hyperlink{#1}{\textcolor{linkc}{slide~\pageref*{#1}}}}
\newcommand{\srefs}[2]{\textcolor{linkc}{slides~\hyperlink{#1}{\pageref*{#1}}--\hyperlink{#2}{\pageref*{#2}}}}
\newcommand{\plainen}[1]{\colorbox{panel}{\parbox{\dimexpr\linewidth-2\fboxsep}{\raggedright\footnotesize\textbf{In plain English:} #1}}\par\vspace{1mm}}
\newcommand{\idcard}[1]{\colorbox{panel}{\parbox{\dimexpr\textwidth-2\fboxsep}{\scriptsize #1}}}
\newcommand{\grp}[2]{\multicolumn{#1}{@{}l}{\textcolor{awsteal}{\textit{#2}}}}
\newcommand{\agi}[4]{\hyperlink{#1}{{\large\bfseries\color{awsteal}#2}\quad{\large\bfseries\color{white}#3}}\\[0.5mm]{\small\color{inkmuted}#4}\par\vspace{4mm}}
\newcommand{\bgframe}[1]{\usebackgroundtemplate{\includegraphics[width=\paperwidth,height=\paperheight]{#1}}}

\arrayrulecolor{rulec}
\lstset{basicstyle=\ttfamily\tiny\color{bodytext},breaklines=true,columns=fullflexible,
  backgroundcolor=\color{panel},frame=single,rulecolor=\color{rulec},
  keywordstyle=\color{awsteal}\bfseries,commentstyle=\color{inkmuted},showstringspaces=false}
\renewcommand{\arraystretch}{1.1}
\usetikzlibrary{arrows.meta,positioning,shapes.geometric,calc,fit,backgrounds,shadows}
\tikzset{box/.style={draw=awsblue,fill=panel,rounded corners=2pt,align=center,font=\scriptsize,inner sep=3pt,text=bodytext},
         arr/.style={-{Stealth[length=2mm]},thick,awsblue}}
\pgfplotsset{every axis/.append style={axis line style={rulec},tick style={rulec},
  tick label style={text=bodytext},label style={text=bodytext}}}

% numbered table / figure labels (one index across the deck)
\newcounter{tabno}\newcounter{figno}
\newcommand{\tabcap}[1]{\refstepcounter{tabno}{\scriptsize\bfseries\color{awsteal}Table~\thetabno\hspace{0.6em}{\color{bodytext}#1}}\par\vspace{0.2mm}}
\newcommand{\figcap}[1]{\refstepcounter{figno}{\scriptsize\bfseries\color{awsteal}Figure~\thefigno\hspace{0.6em}{\color{bodytext}#1}}\par\vspace{0.2mm}}

\newcommand{\deckfooter}{\textcopyright{} \the\year, Amazon Web Services, Inc. or its affiliates.}

% ---- beamer look ---------------------------------------------------
\setbeamercolor{background canvas}{bg=bgdark}
\setbeamercolor{normal text}{fg=bodytext,bg=bgdark}
\setbeamercolor{structure}{fg=awsblue}
\setbeamercolor{frametitle}{fg=white}
\setbeamercolor{itemize item}{fg=awsteal}
\setbeamercolor{itemize subitem}{fg=inkmuted}
\setbeamercolor{enumerate item}{fg=awsteal}
\setbeamercolor{alerted text}{fg=awsmagenta}
\setbeamertemplate{navigation symbols}{}
\setbeamertemplate{itemize items}[square]
\setbeamertemplate{enumerate items}[default]
\setbeamertemplate{caption}[numbered]
\setbeamerfont{itemize/enumerate subbody}{size=\scriptsize}
\setbeamersize{text margin left=8mm,text margin right=8mm}

% section tag shown above each frame title (no divider slides)
\gdef\secname{}
\newcommand{\sect}[1]{\section{#1}\gdef\secname{#1}}
\setbeamertemplate{frametitle}{%
  \nointerlineskip\vskip3.2mm
  \ifx\secname\empty\else{\fontsize{6.5}{7}\selectfont\bfseries\color{sectag}\MakeUppercase{\secname}}\par\vskip0.9mm\fi
  {\fontsize{12.5}{14.5}\selectfont\bfseries\color{white}\insertframetitle}\par\vskip0.6mm}
\setbeamertemplate{headline}{}
\setbeamertemplate{footline}{%
  \hbox to\paperwidth{\hspace*{8mm}{\fontsize{5}{6}\selectfont\color{footc}\deckfooter}\hfill{\fontsize{6}{7}\selectfont\color{footc}\insertframenumber}\hspace*{8mm}}\vskip2.5mm}

% ---- deck macros ---------------------------------------------------
\newcommand{\hl}[1]{\textcolor{awsteal}{\textbf{#1}}}
\newcommand{\src}[1]{\vfill{\tiny\textcolor{footc}{Source: #1}}}
\newcommand{\flabel}[1]{\label{#1}\hypertarget{#1}{}}
% clickable slide references
\newcommand{\sref}[1]{\hyperlink{#1}{\textcolor{linkc}{slide~\pageref*{#1}}}}
\newcommand{\srefs}[2]{\textcolor{linkc}{slides~\hyperlink{#1}{\pageref*{#1}}--\hyperlink{#2}{\pageref*{#2}}}}
\newcommand{\plainen}[1]{\colorbox{panel}{\parbox{\dimexpr\linewidth-2\fboxsep}{\raggedright\footnotesize\textbf{In plain English:} #1}}\par\vspace{1mm}}
\newcommand{\idcard}[1]{\colorbox{panel}{\parbox{\dimexpr\textwidth-2\fboxsep}{\scriptsize #1}}}
\newcommand{\grp}[2]{\multicolumn{#1}{@{}l}{\textcolor{awsteal}{\textit{#2}}}}
\newcommand{\agi}[4]{\hyperlink{#1}{{\large\bfseries\color{awsteal}#2}\quad{\large\bfseries\color{white}#3}}\\[0.5mm]{\small\color{inkmuted}#4}\par\vspace{4mm}}
\newcommand{\bgframe}[1]{\usebackgroundtemplate{\includegraphics[width=\paperwidth,height=\paperheight]{#1}}}

\arrayrulecolor{rulec}
\lstset{basicstyle=\ttfamily\tiny\color{bodytext},breaklines=true,columns=fullflexible,
  backgroundcolor=\color{panel},frame=single,rulecolor=\color{rulec},
  keywordstyle=\color{awsteal}\bfseries,commentstyle=\color{inkmuted},showstringspaces=false}
\renewcommand{\arraystretch}{1.1}
\usetikzlibrary{arrows.meta,positioning,shapes.geometric,calc,fit,backgrounds,shadows}
\tikzset{box/.style={draw=awsblue,fill=panel,rounded corners=2pt,align=center,font=\scriptsize,inner sep=3pt,text=bodytext},
         arr/.style={-{Stealth[length=2mm]},thick,awsblue}}
\pgfplotsset{every axis/.append style={axis line style={rulec},tick style={rulec},
  tick label style={text=bodytext},label style={text=bodytext}}}

% numbered table / figure labels (one index across the deck)
\newcounter{tabno}\newcounter{figno}
\newcommand{\tabcap}[1]{\refstepcounter{tabno}{\scriptsize\bfseries\color{awsteal}Table~\thetabno\hspace{0.6em}{\color{bodytext}#1}}\par\vspace{0.2mm}}
\newcommand{\figcap}[1]{\refstepcounter{figno}{\scriptsize\bfseries\color{awsteal}Figure~\thefigno\hspace{0.6em}{\color{bodytext}#1}}\par\vspace{0.2mm}}

\renewcommand{\deckfooter}{Pre-read for the agentic RL deep-dive. Answers can be short; ``don't know yet'' is a useful answer.}

\tikzset{wf/.style={box,font=\fontsize{7}{8.4}\selectfont,minimum height=10mm,text width=21mm,inner sep=2pt},
         wfa/.style={-{Stealth[length=1.8mm]},thick,awsblue},
         lab/.style={font=\fontsize{5.6}{6.6}\selectfont,text=inkmuted},
         num/.style={circle,fill=awsmagenta,text=white,font=\fontsize{5.5}{6}\selectfont\bfseries,inner sep=0.6pt,minimum size=3.2mm}}

\begin{document}

% =====================================================================
\sect{Agentic RL engagement}
\begin{frame}
\frametitle{The agentic RL workflow}
\vspace{-1mm}
{\small\color{inkmuted}A model learns a multi-step job by trying it in a real environment many times and keeping what scores well.\par}
\vspace{3mm}
\centering
\begin{tikzpicture}[node distance=6mm]
  \node[wf] (task) {\textbf{Task + env}\\instruction, files};
  \node[wf,right=of task] (har) {\textbf{Harness}\\agent loop};
  \node[wf,right=of har] (pol) {\textbf{Policy model}\\served by inference};
  \node[wf,right=of pol] (sbx) {\textbf{Tools / sandbox}\\run the actions};
  \node[wf,right=of sbx,draw=awsteal] (rew) {\textbf{Reward}\\verifier scores it};
  \node[wf,below=11mm of sbx,draw=awspurple] (tr) {\textbf{Trainer}\\RL update};
  \node[wf,below=11mm of rew,draw=awsmagenta] (ev) {\textbf{Evaluation}\\held-out tasks};
  \draw[wfa] (task) -- (har);
  \draw[wfa,<->] (har) -- (pol);
  \draw[wfa,<->] (pol) -- (sbx);
  \draw[wfa] (sbx) -- (rew);
  \draw[wfa] (rew.south) -- ++(0,-4mm) -| node[pos=0.25,below,lab]{trajectory + reward} ([xshift=4mm]tr.north);
  \draw[wfa] (tr.west) -| node[pos=0.25,below,lab]{new weights} (pol.south);
  \draw[wfa,dashed,awsmagenta] (tr.east) -- node[below,lab]{checkpoint} (ev.west);
  \node[num] at (task.north west) {1};
  \node[num] at (har.north west) {2};
  \node[num] at (sbx.north west) {3};
  \node[num] at (rew.north west) {4};
  \node[num] at (tr.north west) {5};
  \node[num] at (ev.north west) {6};
\end{tikzpicture}

\vspace{5mm}
\raggedright
\begin{columns}[T,onlytextwidth]
\begin{column}{0.49\textwidth}\footnotesize
\begin{enumerate}\setlength{\itemsep}{0.6mm}
\item \hl{Pick a task} and start a fresh environment for it.
\item The \hl{harness} runs the agent loop: ask the model, parse its tool call, return the result, repeat.
\item \hl{Tools} act in an isolated \hl{sandbox} (shell, files, APIs, a GPU), so mistakes stay contained.
\end{enumerate}
\end{column}
\begin{column}{0.49\textwidth}\footnotesize
\begin{enumerate}\setcounter{enumi}{3}\setlength{\itemsep}{0.6mm}
\item When the agent stops, a \hl{verifier} checks the end state and gives a reward (e.g.\ hidden tests pass).
\item The \hl{trainer} compares several attempts at the same task and makes the better ones more likely (e.g.\ GRPO).
\item \hl{Evaluation} on held-out tasks shows whether the gain is real, not memorized.
\end{enumerate}
\end{column}
\end{columns}
\end{frame}

% =====================================================================
\begin{frame}
\frametitle{The components, and what each one does}
\vspace{-1mm}
{\small\color{inkmuted}Each row is something you own or we build together. Examples are from our Terminal-Bench + KernelBench demo on AgentCore.\par}
\vspace{2mm}
\fontsize{7}{8.6}\selectfont
\renewcommand{\arraystretch}{1.26}
\begin{tabularx}{\textwidth}{@{}>{\raggedright\arraybackslash}p{27mm}>{\raggedright\arraybackslash}X>{\raggedright\arraybackslash}X@{}}
\toprule
\textcolor{awsteal}{\textbf{Component}} & \textcolor{awsteal}{\textbf{Role}} & \textcolor{awsteal}{\textbf{Example}}\\
\midrule
\textbf{Policy model} & the model being trained; it chooses every action & Qwen3-4B / Qwen3.5-9B, open weights\\
\textbf{Tasks + training data} & what the agent practises on, with a train/eval split & $\sim$200 Terminal-Bench shell tasks and 100 KernelBench GPU problems, plus held-out eval sets\\
\textbf{Task definition + reward} & what ``done'' means, scored automatically for each task type & hidden pytest oracle (1 / partial / 0); KernelBench: correct = 0.5, correct and $\geq$1.1$\times$ faster = 1\\
\textbf{Harness} & the agent loop: prompt, tool-call parsing, turn and token budgets & a mini Claude Code loop (8 turns); the real Claude Code CLI as an option\\
\textbf{Tools + skills} & the actions the agent may take, and how-to guides it can read on demand & \texttt{bash}, \texttt{write\_file}, \texttt{read\_file}; skills such as ``self-check your kernel''\\
\textbf{Sandbox environment} & an isolated, fresh machine per attempt, where the tools run & AgentCore microVM; EC2 for multi-container tasks; g6 GPU instance\\
\textbf{Trainer + inference} & generate attempts fast, then update the weights & slime (Megatron + SGLang), fully-async GRPO on 8$\times$L40S\\
\textbf{Evaluation} & measures the gain on tasks the model never trained on & held-out solve rate every 10 steps, plus cost and latency per task\\
\bottomrule
\end{tabularx}
\end{frame}

% =====================================================================
\begin{frame}
\frametitle{Questions to answer before our deep-dive}
\vspace{-1mm}
\fontsize{6.2}{7.3}\selectfont
\renewcommand{\arraystretch}{1.12}
\begin{tabularx}{\textwidth}{@{}>{\raggedright\arraybackslash}p{24mm}>{\raggedright\arraybackslash}X>{\raggedright\arraybackslash}p{44mm}@{}}
\toprule
\textcolor{awsteal}{\textbf{Area}} & \textcolor{awsteal}{\textbf{Key questions}} & \textcolor{awsteal}{\textbf{Why we ask}}\\
\midrule
\textbf{1\; Goal and success} & What should the agent do better than today? Which KPI does it move, and what number counts as success? What fails with today's approach (prompting, SFT, an API model)? & RL optimizes exactly what we measure, so we must agree on it first.\\
\textbf{2\; Model} & Base model and size? Open weights or API-only? Licence limits on fine-tuning? Starting checkpoint? Serving target (latency, cost, hardware)? & RL needs the weights; model size sets GPU count and step time.\\
\textbf{3\; Tasks and data} & Which task types? How many tasks of each type, and where do they come from (logs, a benchmark, synthetic)? Train/eval split? PII or contamination risk? Past trajectories? & Too few or too-easy tasks give no learning signal; the split decides what eval proves.\\
\textbf{4\; Task + reward per type} & For each task type: input, expected output, verifier (tests, exact match, a checker, an LLM judge), reward (binary or partial)? Does it run automatically, in seconds? How could it be gamed? & The reward is the whole training signal. If it can be gamed, it will be.\\
\textbf{5\; Harness, tools, skills} & Which harness (Claude Code, Codex, LangGraph, your own)? Tool APIs with their schemas: read-only or side-effecting, rate limits? Existing skills or playbooks? Turn and token budget? & Train in the same loop the model will be deployed in, or the gains don't transfer.\\
\textbf{6\; Sandbox environment} & What does one episode need: OS/arch, image size, CPU or GPU, network, multi-container, secrets? Episode length? How many in parallel? Must state reset between episodes? & This decides the sandbox type (microVM, EC2, GPU) and the cost per episode.\\
\textbf{7\; Evaluation} & Metrics (solve rate, pass@k, cost and latency per task, safety)? Held-out set or public benchmark? Baseline to beat? Human or LLM-judge spot checks? & Without a fixed eval and a baseline, a better training curve proves nothing.\\
\textbf{8\; Constraints} & Compute budget, AWS region/account? Data residency, security and legal reviews? Timeline? Owners for data, environment, reward and eval? & This sets the scope, and who unblocks each piece.\\
\bottomrule
\end{tabularx}
\vspace{1.5mm}
\colorbox{awsblue!35!panel}{\parbox{\dimexpr\textwidth-2\fboxsep-1pt}{\fontsize{7}{8.4}\selectfont\color{white}%
\textbf{Please bring:} 5 example tasks, each with its expected output and how you grade it today.}}
\end{frame}
\end{document}
